% Variante : anglais, corrigés reportés en fin de partie.
% Même corps que les autres variantes ; seule la ligne \usepackage change.
\documentclass[11pt]{ocots-book}
\usepackage[lang=en, theme=ocots, solutions=end, math=analysis]{ocots}
\lstset{language=[LaTeX]TeX}

\title{Variante : anglais, corrigés reportés}
\author{Prénom \textsc{Nom}}
\date{\today}

\newcommand{\ocotsvariantnote}[1]{\begin{quote}\itshape Variante : #1\end{quote}}

\begin{document}

\begin{lstlisting}
\documentclass[11pt]{ocots-book}
\usepackage[lang=en, theme=ocots, solutions=end, math=analysis]{ocots}
\end{lstlisting}
\ocotsvariantnote{\texttt{lang=en} (le défaut du template est \texttt{lang=fr})
: les intitulés (« Theorem » au lieu de « Théorème »...) et la typographie
suivent l'anglais. \texttt{solutions=end} est ici la même valeur que le
défaut — seule la langue change vraiment par rapport au réglage usuel.}

\ocotscollectsolutions
% Corps partagé par toutes les variantes : aucune option, aucun préambule.
%
% C'est la démonstration de l'étape 2 : le même contenu, sans une ligne de
% différence, se compose en français ou en anglais, avec ou sans corrigés, en
% couleur ou en noir et blanc. Seule la ligne \usepackage change.

\chapter{Chapitre de démonstration}
\begin{chapterintro}
    Une introduction de chapitre, avec \ie un exemple de macro localisée et
    \enquote{un texte cité}.
\end{chapterintro}
\section{Une section}

\begin{theorem}{Un théorème}{demo}
    Énoncé.
\end{theorem}

\begin{proof}
    Démonstration.
\end{proof}

\begin{definition}{Une définition}{}
    Énoncé.
\end{definition}

\begin{remark}
    Une remarque.
\end{remark}

\begin{assumption}
    Une hypothèse, étiquetée automatiquement.
\end{assumption}

\begin{exercise}[points=4]
    Énoncé de l'exercice.
\solution
    Corrigé.
\end{exercise}

\ocotsprintsolutions
\end{document}
